%%%%%%%%%%%%%%%%%%%%%%%%%%%%%%%%%%%%%%%%
%        Author: Bernard Lampe
%   Binary reverse engineering lecture material
%%%%%%%%%%%%%%%%%%%%%%%%%%%%%%%%%%%%%%%%

% import template
\documentclass[12pt]{Training}

% import packages
\usepackage[utf8]{inputenc}
\usepackage[T1]{fontenc}
\usepackage{mathpazo}
\usepackage{graphicx}
\usepackage{booktabs}
\usepackage{listings}
\usepackage{enumerate}
\usepackage{xcolor}

%-------------------------------------------------------------------------------

% set document info
\title{Binary Reverse Engineering Lecture Material}
\author{Dr. Bernard Lampe}
\class{Introduction to Vulnerability Research}
\date{October 31st, 2019}

%-------------------------------------------------------------------------------

% set code listing style
\lstset{
    basicstyle=\ttfamily\small,
    numbers=left,
    tabsize=2,
    keywordstyle=\color{blue}\ttfamily,
    stringstyle=\color{red}\ttfamily,
    commentstyle=\color{green}\ttfamily,
    morecomment=[l][\color{magenta}]{\#}
}

%-------------------------------------------------------------------------------

\begin{document}

% print title
\maketitle

%-------------------------------------------------------------------------------

\section*{Introduction}
\begin{enumerate}
    \item Why reverse?
    \begin{enumerate}
        \item Source code is often unavailable: firmware, games, malware, proprietary protocol implementations, vendor drivers.
        \item To understand execution at the machine level, which respects no abstraction: the CPU executes what the compiler emitted, not what the developer intended.
        \item To verify source claims when auditing: compilers transform code, and the shipped binary is the ground truth (optimizations, inlining, dead code elimination).
        \item To recover structure: symbol stripping, obfuscation, and packing commonly accompany anything valuable to protect.
    \end{enumerate}
    \item When would you reverse (give example situations)?
    \begin{enumerate}
        \item Fuzzing an undocumented file format: understanding the parser directs the fuzzer and seeds the corpus.
        \item Vulnerability research on closed source targets (browsers, drivers, IoT firmware) where a large part of the attack surface is compiled code.
        \item Building a sanitizer/instrumentation harness: you need to know where the input crosses the trust boundary before you can hook it.
        \item Writing an exploit: kernel structures, driver interaction, and mitigation layout are read from the binary and from the running system.
        \item Malware analysis: capability identification, C2 protocol recovery, IOC production.
        \item Interoperability: reimplementing a protocol when the vendor refuses to document it.
    \end{enumerate}
    \item How does reversing fit into the larger VR picture?
    \begin{enumerate}
        \item It is the static-analysis half of vulnerability research when no source exists; it fills the same role source auditing does for open targets: static analysis, with high false-positive cost and deep target knowledge required, when the source is not available.
        \item It feeds dynamic analysis: reversing produces breakpoints, input grammars, harness entry points, and expected invariants for debugging and fuzzing.
        \item It is iterative with exploitation: superseding assumptions (allocator layout, mitigation status) requires reading the shipped code.
    \end{enumerate}
    \item What are the legal aspects of reversing?
    \begin{enumerate}
        \item This is not legal advice. Laws differ by jurisdiction (US DMCA, EU Software Directive, etc) and are applied unevenly.
        \item Interoperability and research exceptions generally exist (e.g. reverse engineering to achieve interoperability in the EU; fair-use argumentation in the US, cf. Sega v. Accolade, Sony v. Connectix), but anti-circumvention provisions and EULAs complicate them.
        \item Practical rules: reverse what you own or are licensed to run, do not redistribute their code or derived assets, do not attack third parties from your findings, and get written authorization for anything that touches systems you do not own.
    \end{enumerate}
\end{enumerate}

\section*{Overview of How Compilers Work}
\begin{enumerate}
    \item Go from C to binary and back, to get an idea of what happens under the hood.
    \begin{enumerate}
        \item A C source file is preprocessed (macro expansion, header inclusion, conditional compilation), then lowered through successively less abstract IRs (in LLVM: Clang AST $\to$ LLVM IR $\to$ selection DAG $\to$ MachineInstr $\to$ MCInst) and encoded into object code.
        \item The reverse path, taken by decompilers, loses information at every lowering step. Names, types, comments, control-flow intent are gone. This is why decompilation is lossy and why guided reversing is needed.
        \item Anything the compiler knows from the source (struct layout, constant expressions) may be folded into immediates; recovery ones are inferred by the decompiler as best-effort.
    \end{enumerate}
    \item Front end of the compiler (lexical analysis, syntax analysis per a grammar, semantic analysis/type checking)
    \begin{enumerate}
        \item Lexical analysis: characters to tokens.
        \item Syntax analysis: tokens to a parse tree per a grammar.
        \item Semantic analysis: type checking, scope resolution, implicit conversions.
        \item Front end output (AST) is lowered to IR; the IR platform-independence is what makes retargetable code generation possible.
    \end{enumerate}
    \item Back end of the compiler (optimization, register spilling, code generation)
    \begin{enumerate}
        \item Optimization at IR level: dead code elimination, constant propagation, inlining, loop unrolling, vectorization.
        \item Instruction selection: IR to machine instructions; register allocation: instructions to physical registers with spilling to stack slots when registers run out.
        \item Code generation and emission: scheduling, peephole, object file writing (relocations, symbol tables).
        \item Implication for reversing: optimization changes code shape (inlined functions disappear, spilled values are untracked). When possible compare -O0 through the shipped optimization level to classify structures.
        \item Idioms that identify an optimization at a glance: a \lstinline{jmp} over an unreachable block (dead code), a call whose result is stored but never read (eliminated), a loop replaced by a closed-form expression (strength reduction, loop unrolling), and a \lstinline{memset}/\lstinline{memcpy} call where the source had a scalar loop. Recognizing these is most of reading optimized output.
        \item \lstinline{-fomit-frame-pointer} removes the frame-pointer chain that makes a stack walk trivial; the unwind tables (\lstinline{.eh_frame}, \lstinline{.pdata}) are then the only reliable call-graph ground truth. Keep the tables even when the code is stripped.
        \item Debug info is a separate, optional artifact: DWARF on Linux/macOS (\lstinline{-g}, \lstinline{.debug_*} sections), PDB on Windows. \lstinline{objcopy --only-keep-debug} splits it from the shipped binary, and the split is itself a reversing shortcut when it leaks.
        \item Intermediate representation is readable: \lstinline{clang -S -emit-llvm} prints LLVM IR, which is far closer to the source than assembly. When a vendor ships IR or bitcode (firmware, JIT caches, Swift modules), read that before the machine code.
\end{enumerate}
\end{enumerate}

\section*{Object Formats and Linking}
\begin{enumerate}
    \item Before code is data: object formats. ELF is the Linux/Unix standard, Mach-O the macOS/iOS one, PE the Windows one.
    \item Sections: executable code (.text/\_\_TEXT,\_\_text), initialized data (.data/\_\_DATA,\_\_data), uninitialized data (.bss), read-only data (.rodata/\_\_TEXT,\_\_cstring), unwind/debug info.
    \item Symbol tables map names to addresses; stripping removes them but the code remains. Relocations are patch points applied at link/load time; resolving them is how position-independent calls are walked.
    \item Import tables (PLT/GOT on ELF, stubs and lazy binding on Mach-O) show which library functions the code calls without reversing a byte of it: usually the first stop when triaging a binary.
    \item File formats are themselves parsers: malformed headers and section tables crash tools (disassemblers and loaders are attack surface too).
\end{enumerate}

\textbf{ELF, the Linux/macOS-unix default}
\begin{enumerate}
    \item Read the header, then the program headers, then the section headers: \lstinline{e_ident} identifies the class (32/64) and endianness, \lstinline{e_type} distinguishes executable from shared object from core, and the program headers are what the loader actually consumes.
    \item The dynamic section is the linking contract: \lstinline{DT_NEEDED} lists the libraries, \lstinline{DT_SYMTAB}/\lstinline{DT_STRTAB}/\lstinline{DT_GNU_HASH} locate the symbol and hash tables, \lstinline{DT_INIT}/\lstinline{DT_INIT_ARRAY} name the constructors that run before \lstinline{main}, and \lstinline{DT_FLAGS} records \lstinline{NOW} binding.
    \item \lstinline{PT_GNU_RELRO} marks the range the loader re-protects read-only after relocation; \lstinline{PT_GNU_STACK} encodes the executable-stack request, and \lstinline{PT_LOAD} entries carry the actual permissions. These are the on-disk mirror of the mitigation map: RELRO is visible in \lstinline{PT_GNU_RELRO}, the executable-stack request in \lstinline{PT_GNU_STACK}, NX in the \lstinline{PT_LOAD} permissions, and the loader's symbol protection in full RELRO.
    \item The PLT/GOT pair is the import mechanism: each unresolved call goes through a stub that pushes a relocation index and jumps to the resolver; after the first call the slot holds the resolved address (lazy binding). \lstinline{LD_BIND_NOW} or \lstinline{-z now} forces eager resolution, and full RELRO re-protects \lstinline{.got} afterwards --- the hardening that closes the GOT-overwrite primitive.
    \item \lstinline{STT_GNU_IFUNC} symbols (glibc's \lstinline{memcpy}, \lstinline{strlen} selection) are resolved at load time by running code, so a resolver is an execution path that exists before \lstinline{main}. Rare, but a real reversing and auditing surprise.
    \item \lstinline{PT_TLS} and the thread-pointer register (\lstinline{fs}/\lstinline{gs} on x86, \lstinline{tpidr_el0} on ARM64) locate thread-local storage; a wrong \lstinline{tpidr} recovery is a common decompiler failure on Android.
    \item Stripping removes \lstinline{.symtab}, not the dynamic symbol table: \lstinline{.dynsym} still names every imported and exported function, so imports survive stripping and are the first thing to read. \lstinline{--only-keep-debug} plus \lstinline{--strip-debug} produces the common shipped split (debug file kept by the vendor).
    \item Unwind data (\lstinline{.eh_frame}, \lstinline{.eh_frame_hdr}, \lstinline{.ARM.exidx}) is present even in stripped binaries that were built with exceptions; it is a function-boundary oracle the stripper usually forgets to remove.
\end{enumerate}

\textbf{Mach-O, the macOS/iOS default}
\begin{enumerate}
    \item Load commands (\lstinline{LC_SEGMENT_64}, \lstinline{LC_MAIN}, \lstinline{LC_DYLD_INFO}) replace ELF's program headers; segments are named \lstinline{__TEXT}, \lstinline{__DATA}, \lstinline{__LINKEDIT}.
    \item Imports go through \lstinline{__stubs}/\lstinline{__stub_helper} (the PLT/GOT analogue); modern binaries replace them with chained fixups (\lstinline{LC_DYLD_CHAINED_FIXUPS}), where each pointer carries its own target encoding.
    \item Objective-C and Swift metadata (\lstinline{__objc_classlist}, \lstinline{__objc_methname}, \lstinline{__swift5_*}) is high-level structure the compiler emitted into the binary: method names, class layouts, and protocol conformances survive stripping, which makes a stripped Objective-C binary far easier to read than a stripped C one.
    \item On iOS the shared cache (\lstinline{dyld_shared_cache}) holds the system libraries pre-linked; reversing it requires the \lstinline{dyld_shared_cache_util} extraction step.
\end{enumerate}

\textbf{PE, the Windows default}
\begin{enumerate}
    \item The DOS header's \lstinline{e_lfanew} points at the NT headers; the optional header holds the entry point (\lstinline{AddressOfEntryPoint}), image base, section alignment, and the data directories.
    \item Imports are described by \lstinline{IMAGE_IMPORT_DESCRIPTOR} entries naming DLLs and their \lstinline{IMAGE_THUNK_DATA} arrays, which are the IAT slots the loader patches. Delayed and ordinal imports are the variants to recognize.
    \item Exports (\lstinline{IMAGE_EXPORT_DIRECTORY}) name the DLL's own functions: on a shipped Windows binary, exports are the most reliable symbol source.
    \item \lstinline{RUNTIME_FUNCTION} entries (\lstinline{.pdata}) are the x64 exception/unwind table; they are machine-readable function boundaries and therefore a reversing aid even without symbols.
\end{enumerate}

\section*{Tools}
\begin{enumerate}
    \item Reversing tools like Radare, IDA, Binja, objdump, etc
    \begin{enumerate}
        \item Disassembly + decompilation: IDA Pro (Hex-Rays, the industry standard; commercial), Binary Ninja (newer, strong ILs for scripted analysis), Ghidra (NSA open source, includes decompiler), radare2/rizin (free, scriptable, steep learning curve), Hopper (macOS).
        \item Command line for quick looks: objdump (binary sizes, sections, symbols, simple disasm), llvm-objdump, otool/nm/lipo on macOS, readelf on ELF.
        \item Dynamic analysis: gdb (gdb-multiarch recommended for remote targets), LLDB, debuggers built into the above; Frida for scriptable hooking of live processes; ltrace/strace/turbostat for syscall-level views.
        \item Scriptability is a core requirement: real reversing is automated (batch findings across thousands of functions; write plugins or IDAPython/gdb scripts).
        \item Library and framework choices: \lstinline{LIEF} and \lstinline{pefile} for parsing formats in scripts, \lstinline{Capstone}/\lstinline{Zydis} for disassembly, \lstinline{Keystone} for assembling, \lstinline{Unicorn}/\lstinline{Qiling} for emulation, and \lstinline{angr}/\lstinline{miasm}/\lstinline{Triton} for symbolic work. These are the primitives a custom reversing tool is built from.
        \item Symbol servers and signature databases (Microsoft's public symbol server, \lstinline{FLIRT} signatures, \lstinline{Lumina}, \lstinline{FunctionID}) turn a stripped binary back into a named one. Reading \lstinline{pdb} paths out of a PE or \lstinline{.gnu_debuglink} out of an ELF is the first move of a triage.
        \item Debugger integration is what separates a disassembler from a reversing tool: the debugger is the ground truth, and a decompiler that cannot be checked against a run is a hypothesis generator.
    \end{enumerate}
\end{enumerate}

%----------------------------------------------------------------------------------------

\section*{Anti-Reversing Techniques}
\begin{enumerate}
    \item Families:
    \begin{enumerate}
        \item \textbf{Packing}: UPX, VMProtect/Themida style virtualizers, custom crypters. Detection: section entropy high in exactly one section, few imports, self-modifying memory at runtime. Defeat: run and dump (user-mode unpackers, or break at the original entry point in a debugger), or emulate.
        \item \textbf{Anti-debugging}: \lstinline{IsDebuggerPresent}, PEB/proc status checks, \lstinline{ptrace(PT_DENY_ATTACH)} on macOS, timing checks, exception filters as control flow. Defeat: patch out the calls once located, or attach before they run.
        \item \textbf{Obfuscation}: control-flow flattening (dispatcher + state variable shape in the decompiler), opaque predicates, instruction substitution, string encryption. Defeat: patience + scripting; flattening complicates control flow, but the data flow remains traceable.
        \item \textbf{Integrity checks}: self-hash of code sections at runtime. Defeat: patch + recompute/resign, or break the checker call site.
    \end{enumerate}
    \item The cost logic: anti-reversing raises the attacker's analysis effort per function, not the bug count. When a target is heavily obfuscated, put scripting/automation effort in first --- manual reading rarely amortizes.
    \item Concrete detection tells for the families above:
    \begin{enumerate}
        \item \textbf{Packing}: section entropy near 8.0 in exactly one section, a section name that is not standard (\lstinline{UPX0}/\lstinline{UPX1}), few or no imports, an entry point outside the first section, a writable-and-executable section at load, and a tiny on-disk-to-in-memory size ratio.
        \item \textbf{Anti-debugging}: imports of \lstinline{IsDebuggerPresent}, \lstinline{CheckRemoteDebuggerPresent}, \lstinline{NtQueryInformationProcess} with \lstinline{ProcessDebugPort}/\lstinline{ProcessDebugFlags}/\lstinline{ProcessDebugObjectHandle}, \lstinline{NtSetInformationThread} with \lstinline{ThreadHideFromDebugger}, \lstinline{OutputDebugString} with a last-error check, \lstinline{ptrace(PT_DENY_ATTACH)} on macOS, a \lstinline{rdtsc}/\lstinline{clock_gettime} pair, and \lstinline{int 2d}/\lstinline{int 3} or \lstinline{popf} used as a check.
        \item \textbf{Obfuscation}: control-flow flattening (one dispatcher block with a state variable compared in every successor), MBA expressions, opaque predicates that are always true but not provably so, encrypted strings decrypted at use, and instruction substitution. LLVM-based obfuscators (\lstinline{OLLVM}, \lstinline{tigress}) leave recognizable shapes.
        \item \textbf{Integrity checks}: a hash or CRC over the code section, computed and compared at startup; the comparison site is the patch point, and the recompute is the resign step.
    \end{enumerate}
    \item Defeat toolkit, in cost order: patch the check once located; hide from the check (\lstinline{ScyllaHide}, \lstinline{TitanHide}); unpack to the original entry point (\lstinline{OEP}) and dump; emulate (\lstinline{Unicorn}/\lstinline{Qiling}) so the check never runs; and symbolically simplify (\lstinline{Triton}/\lstinline{miasm}) when the obfuscation is arithmetic. Every one of these is itself a tooling project, which is the effort the protection demands.
    \item The honest framing for the report: anti-reversing is not a bug and not a mitigation. It is a burden on the \emph{auditor}, and the correct response is a time/effort estimate, exactly as a mitigation warrants an estimate of what its bypass requires.
\end{enumerate}

\section*{Architectures x86, x86\_64, ARM, ARM64, MIPS, PPC}
\begin{enumerate}
    \item Things about architectures.
    \begin{enumerate}
        \item Instruction sets come in orthogonal and non-orthogonal flavors, fixed and variable size encodings, load-store vs. memory-register models. Table:
        \begin{center}
        \footnotesize
        \begin{tabular}{@{}lll@{}}
        \toprule
        Arch & Encoding & Notes \\
        \midrule
        x86 & variable (1--15 bytes) & dense, memory-register ops, hard to decode \\
        x86\_64 & variable & +8 GPRs, RIP-relative addressing \\
        ARM & fixed 4 bytes (A32) & condition codes on every instruction \\
        Thumb2 & 2/4 byte mixture & ARMv7+; mode ambiguity in disassembly \\
        ARM64 & fixed 4 bytes & 31 GPRs, no condition flags carry-over \\
        MIPS & fixed 4 bytes & load-store, delay slots, canonical in embedded \\
        PPC & fixed 4 bytes & load-store, big endian in older systems \\
        \bottomrule
        \end{tabular}
        \end{center}
        \item Calling conventions: how arguments are placed, where the return value is returned, what registers are callee-saved, how structs/classes are constructed, all differ per ABI.
        \begin{enumerate}
            \item System V AMD64 (Linux/macOS): args in rdi, rsi, rdx, rcx, r8, r9; return in rax; rbx, rbp, r12--r15 callee-saved; return address and frame pointer on the stack.
            \item AArch64: args in x0--x7; return in x0; x19--x28 and x29 (FP) and x30 (LR) callee-saved; LR in a register (implicit push of x30 in prologue).
            \item ARM32: r0--r3 args; LR in r14; condition meaningful for Thumb/ARM mode selection (LSB of function pointers indicating Thumb).
            \item Knowledge of the convention where the target was compiled is a prerequisite to reading prologues, argument registers, and call sites.
        \end{enumerate}
        \item Features that affect reversing: endianness (bi-endian ARM/MIPS/PPC), byte order of embedded filesystem dumps (big endian MIPS/PPC still common in embedded), SIMD/vector extensions (NEON, SSE/AVX, AltiVec) changing data-flow shape, thus decompiler capability.
        \item ARM64 hardening shapes are visible in the prologue and epilogue and change what a stack overflow grants: \lstinline{paciasp}/\lstinline{autiasp} (pointer authentication on the saved \lstinline{x29}/\lstinline{x30}), \lstinline{bti c} (branch-target identification, the ARM64 CFI landing pad), and \lstinline{irg}/\lstinline{stg} (memory tagging). A reversing session that does not name these will mis-read the prologue.
        \item x86 hardening shapes: \lstinline{endbr64} at every indirect-call target (CET, the counterpart of \lstinline{bti c}), shadow-stack instructions, and \lstinline{movaps} alignment requirements that make a pivoted stack fault before the chain runs.
        \item ABI details that recur: x86-64 SysV passes a variadic call's vector-register count in \lstinline{al}; a struct return uses a hidden pointer in the first argument register; ARM64 returns small structs in registers, large ones through \lstinline{x8}. A decompiler that guesses the convention wrong produces a plausible but false signature.
        \item Flavour variants worth knowing: ARM EABI hard-float versus soft-float (argument registers versus a software convention), MIPS \lstinline{o32}/\lstinline{n32}/\lstinline{n64}, and x32. Each changes the argument registers and therefore every recovered prototype.
        \item The audit consequence: an architecture the auditor cannot read is an attack surface the auditor cannot assess. The same bug in a driver on MIPS (delay slots, unaligned access) is a different bug than on x86-64, which is why vendor drivers are the place to apply this section.
    \end{enumerate}
\end{enumerate}

\section*{Disassemblers}
\begin{enumerate}
    \item How disassemblers are written.
    \begin{enumerate}
        \item Ability to decode all instruction encodings of a target is the easy part: table-driven decoders map opcodes to mnemonic/operand tuples. The hard part is control: what bytes are instructions and where functions begin.
        \item Linear sweep: decode from the start of the section to the end. Simple, breaks on data-in-code, on padding, on jump tables that contain addresses, and on relocation stubs.
        \item Recursive descent: start at known entry points, follow control flow (branches), avoid decoding outside reachable ranges. Misses indirect jumps/calls whose target is computed; hard cases require value-set analysis or speculation.
        \item Symbol recovery: from symbol tables when present, from imported function signatures otherwise (PLT resolution), from calling conventions and known library signatures (FLIRT-style signature matching) otherwise.
        \item Ambiguity in Thumb/ARM mode switching, x86 variable length instruction overlap, and embedded padding are ordinary disassembler failure modes: treat disassembler output with some skepticism.
        \item Symbol recovery, in order of preference: the symbol table when present; the dynamic symbol table when only imports/exports survive; import stubs and their call sites (PLT/GOT, Mach-O stubs, PE IAT); library signatures and pattern databases (\lstinline{FLIRT}, \lstinline{Lumina}, \lstinline{FunctionID}) when nothing else remains; and, last, function-boundary heuristics (prologue patterns, alignment, \lstinline{noreturn} calls, unwind tables).
        \item The \lstinline{noreturn} heuristic is why a call to \lstinline{abort} or \lstinline{exit} ends a function without a return instruction; the counterpart is a tail call, which begins with a jump and no prologue. Both are common false function boundaries.
        \item Jump tables are the recurring data-in-code case: a \lstinline{switch} compiles to a table of addresses loaded through a scaled index, and a linear sweep will decode the table as instructions. Recursive descent avoids the table but must know that the indirect jump is a jump.
        \item x86 instruction overlap is the linear sweep's second failure mode: a byte sequence can decode two ways depending on the start offset, so a single wrong byte misaligns the whole rest of the section.
        \item Disassembler engines are libraries now (\lstinline{Capstone}, \lstinline{Zydis}, \lstinline{XED}); building a custom tool means choosing one and supplying the boundary heuristics, not writing the decoder.
    \end{enumerate}
\end{enumerate}

\section*{Decompilers}
\begin{enumerate}
    \item How decompilations are not unique.
    \begin{enumerate}
        \item A decompiler recovers a program consistent with the binary semantics; many source-level programs are consistent, and there is no canonical reverse mapping.
        \item Redeclaration (many C programs compile to the same binary), undefined structure layout, and unreachable code are the largest classes of ambiguity.
        \item The output is a hypothesis, not truth; when a decompiler and a debugger disagree, the debugger (or a run) is right.
    \end{enumerate}
    \item Pitfalls of decompilers, like function inlining and optimizations
    \begin{enumerate}
        \item Inlining erases the call edge: two functions appear as one, types are mixed, the decompiler's inferred prototype is wrong if the call disappeared.
        \item Switch tables are expanded to chains of compares, arrays fold to loops with unrolled bodies: readable structure is lost.
        \item Optimization is a transformation, not a document: branch folding produces single-trace functions that appear to have no condition where the source had many.
        \item Heap/stack register usage is a large area of decompiler weakness: fastcall setup, untyped pointer casts, and goto-based state machines all produce wrong-looking output.
        \item Trust hierarchy: hand disassembly $>$ static decompiler $>$ (both wrong when the binary is packed, self-modifying, or has just unknown encodings).
        \item How a decompiler works, in one line: lift each instruction into a small IR, build the control-flow graph, run SSA construction, then run the classic compiler optimizations backwards (copy propagation, dead-code elimination, expression folding, structure recovery) to print C. Every step is an inference, and each has a known failure mode.
        \item Recurring failure modes to expect: \lstinline{setjmp}/\lstinline{longjmp} and exception edges (the CFG loses the second path), \lstinline{alloca} and variable-length arrays (a dynamic stack frame), inline assembly (opaque to the lifter), tail calls (a missing call edge), \lstinline{varargs} (a prototype that hides the real arguments), and \lstinline{__int128}/floating point (operations the IR may not model exactly).
        \item Structure and type recovery is the weakest part: a decompiler invents a struct where the source had an array of scalars, or an \lstinline{int} where the source had a pointer. Treat every recovered type as a hypothesis to be checked against a run.
        \item Decompiler output is still the best starting document for a large binary: it is searchable, annotatable, and re-nameable. The discipline is to record the confidence of each recovered name and type rather than to trust the printed C.
    \end{enumerate}
\end{enumerate}

\section*{Approach: How to Structure a Reversing Session}
\begin{enumerate}
    \item Triage first: \lstinline{file}, \lstinline{otool -L}/\lstinline{readelf -d}, strings, packer/entropy check. What is it, for what platform, packed or not?
    \item Map the surface: imports, exported functions, entry points, and main loops. Locate input processing (syscalls, file/network reads) since that is where attackable behavior lives.
    \item Automate what you can: scripting disassemblers to bucket functions by size, complexity, and xref count focuses human attention (the auditor's time is the scarce resource, so the tool's job is to rank what to read).
    \item Dynamic: attach a debugger, provide real input, and observe execution; the interplay of static and dynamic is what makes reversing fast.
    \item Document as you go: rename functions, label structures, leave notes in your tool. Reversing is cumulative work; a well-labeled IDB is the deliverable.
    \item The triage steps in order, with the tool for each: \lstinline{file} and the magic bytes; \lstinline{binwalk} for firmware and container images; an entropy plot and section table for packing; \lstinline{strings}/\lstinline{FLOSS} for strings (including obfuscated ones); a YARA scan for known malware and components; and imports plus the entry point for the program's shape.
    \item When the static picture stalls (packed, self-modifying, encrypted strings), go dynamic: run it under a debugger or an emulator (\lstinline{Qiling}/\lstinline{Unicorn}/\lstinline{speakeasy}), let it unpack itself, then dump the memory image and reverse the dump. This is the reversing analogue of skipping the outer layer: feed the already-unpacked image to the analysis instead of fighting the packer.
    \item The reversing deliverable is a written understanding, not a screenshot of a decompiler: the call graph, the input parser's structure, the object layouts, and the invariants the code assumes. That document is what the exploit work and the report consume.
    \item Reversing is iterative with exploitation: a reversing session produces the harness entry points, the breakpoints and expected invariants for dynamic analysis, and the mitigation and layout facts for exploit construction.
\end{enumerate}

\section*{Reversing Examples}
\begin{enumerate}
    \item Example x86, x86\_64, ARM, ARM64, MIPS and PPC. Compile a non-trivial (but smaller) program and compare architectures; the lab document walks a live session on this module's challenge binary (\lstinline{binary_re/targets/checks.c}).
    \begin{enumerate}
        \item The same C source, compiled at -O0 for two architectures (clang on macOS: \lstinline{clang -arch x86_64 -arch arm64 -c hello.c}):
        \footnotesize
        \begin{center}
        \begin{tabular}{@{}ll@{}}
        \toprule
        x86\_64 & ARM64 \\
        \midrule
        \lstinline|pushq %rbp| & \lstinline|sub sp, sp, #0x20| \\
        \lstinline|movq %rsp, %rbp| & \lstinline|stp x29, x30, [sp, #0x10]| \\
        \lstinline|subq $0x10, %rsp| & \lstinline|add x29, sp, #0x10| \\
        \lstinline|leaq "..."(%rip), %rdi| & \lstinline|adrp x0, l_.str@PAGE; add x0, x0, ...| \\
        \lstinline|callq _printf| & \lstinline|bl _printf| \\
        \lstinline|xorl %eax, %eax| & \lstinline|ldr w0, [sp, #0x8]| \\
        \lstinline|popq %rbp; retq| & \lstinline|ldp x29, x30, [sp, #0x10]; add sp, #0x20; ret| \\
        \end{tabular}
        \end{center}
        \item Read it: x86 stores the frame pointer on the stack at every call (the return address is a stack pop away), ARM64 keeps the return address in x30 (LR) and must save it in the prologue --- the \lstinline{stp}/\lstinline{ldp} pair writes/loads both callee-saved registers in one instruction.
        \item Note the ARM64 pattern \lstinline{adrp/add}: PC-relative 4K-page addressing needs two instructions to form a 64-bit address; on x86\_64 a single RIP-relative \lstinline{lea} handles it. Recognizing these idiom pairs is half of reading new architectures.
        \item Follow the call convention: \lstinline{bl} writes the return address into x30 (LR); \lstinline{ret} jumps via x30. That is why overwriting the saved x30 slot in a stack buffer overflow redirects execution on ARM64 the same way overwriting the saved return address does on x86.
    \end{enumerate}
    \item Real-world examples.
    \begin{enumerate}
        \item Take a known CVE and its patch, and reproduce the bug function from the pre-patch binary; identify the missing bounds check in the disassembly. Good practice targets with published writeups: OpenSSL CVE-2014-0160 (Heartbleed), Stagefright libstagefright.so (CVE-2015-3864, integer overflow in the MP4 parser), AppWeb CVE-2018-8715 or older auth-bypass issues.
        \item Exercises: find the comparisons guarding the vulnerable path (the bound check that is missing or wrong), the call into memcpy-like routines, and the length argument. Diff patched and unpatched versions side by side with a binary diff tool (BinDiff, Diaphora) to see the fix as code motion.
    \end{enumerate}
    \item Perhaps tear down an embedded system and then show the following reversing.
    \begin{enumerate}
        \item Embedded teardown flow: obtain firmware (vendor download, SPI flash dump, UART shell), unpack (\lstinline{binwalk} for filesystem images: SquashFS, JFFS2, CramFS), identify binaries (busybox plus vendor daemons), and audit the vendor-written parts (custom code is where the bugs are).
        \item Common bug classes found this way: hardcoded credentials, unauthenticated update acceptance, command injection through configuration values, stack overflows in command parsers.
    \end{enumerate}
    \item Show example of an application process doing network stuff.
    \begin{enumerate}
        \item Reverse triage of a network daemon: find \lstinline{socket}/\lstinline{bind}/\lstinline{accept} (or higher wrappers), follow the buffer that \lstinline{recv} fills into parsing functions, and identify length fields copied through.
        \item Using a debugger: break on the parse entry with a malformed client packet; watch parameter registers under the target's calling convention. This is the binary-side analogue of auditing a web server's header parser with source.
    \end{enumerate}
    \item Show example of reversing the bootloader to setup a chip.
    \begin{enumerate}
        \item Boot chains: mask ROM boot (immutable, vendor loaded) verifies and loads the bootloader (U-Boot, LittleKernel, vendor ones) which verifies and loads the kernel --- verification done by signature check or hash against e-fuses/OTP areas of the chip.
        \item Reversing the bootloader usually means: find the signature verification function, identify what it compares (and whether the comparison is early-exit or constant-time), and establish whether failure paths fall through to by-the-same-code loading. Documented cases (e.g. early Android boot ROM bugs) made this class of bug historically available.
    \end{enumerate}
    \item Other ideas, reverse a proprietary filesystem, proprietary algorithm?
    \begin{enumerate}
        \item Proprietary filesystem: reverse the on-disk format by starting with mount logic and its section parsing; document the superblock and inode layout as C structs in your tool; compare against garbage file-system images you create with partially randomized bytes.
        \item Proprietary algorithm (custom crypto, license checking, game save obfuscation): find the entry by following its caller context (where data is transformed), then instrument inputs/outputs to get pairs of (plaintext, ciphertext) and infer the transformation; static plus dynamic yields the fastest recoveries.
    \end{enumerate}
    \item Reverse a crash from a fuzzing campaign back to source: take a minimized input, break on the parser entry, and read the length and pointer arguments in the registers. The reversing session answers the two questions the sanitizer report only states: which check is missing, and what the attacker controls.
    \begin{enumerate}
        \item The specific exercise is this module's own target: \lstinline{checks.c} in \lstinline{binary_re/targets/} is small enough to reverse by hand and layered enough to teach the triage order (magic, length, checksum, then the bug).
    \end{enumerate}
    \item Reverse a driver's \lstinline{ioctl}/\lstinline{mmap} without source, which is the vendor-driver case: find the \lstinline{file_operations} table through the driver's registration call, then read each verb's argument copy and bounds check. The same audit applies with source in hand: read each verb's argument copy and bounds check in the registration path.
    \item Diff reversing: when a patch exists, diff the patched and unpatched binaries (\lstinline{BinDiff}, \lstinline{Diaphora}, \lstinline{radiff2}) and read the change as code motion. The added check names the bug, which is the fastest root cause available.
    \item Reversing a protocol implementation from a captured session: locate the \lstinline{recv} call site, name the message parser, and rebuild the wire format as a grammar. That grammar is the generative fuzzer's input.
\end{enumerate}

\begin{sectionbox}
\begin{lstlisting}[language={[x86masm]Assembler}, caption={ARM64 slice of \texttt{hello.c} at -O0 (objdump -d)}]
_main:
       0:	ff 83 00 d1	sub	sp, sp, #0x20
       4:	fd 7b 01 a9	stp	x29, x30, [sp, #0x10]
       8:	fd 43 00 91	add	x29, sp, #0x10
       c:	08 00 80 b9	mov	w8, #0x0
      10:	e8 0b 00 b9	str	w8, [sp, #0x8]
      14:	bf c3 1f b8	stur	wzr, [x29, #-0x4]
      18:	00 00 00 90	adrp	x0, l_.str@PAGE
      1c:	00 00 00 91	add	x0, x0, l_.str@PAGEOFF
      20:	00 00 00 94	bl	_printf
      24:	e0 0b 40 b9	ldr	w0, [sp, #0x8]
      28:	fd 7b 41 a9	ldp	x29, x30, [sp, #0x10]
      2c:	ff 83 00 91	add	sp, sp, #0x20
      30:	c0 03 5f d6	ret
\end{lstlisting}
\end{sectionbox}

\section*{Conclusion}
Reversing asks the same analysis questions as source auditing (where does input enter, how is it validated, what objects live where) but without source code. Modern tooling (decompilers, scriptable disassemblers, dynamic instrumentation) has moved the bottleneck from reading assembly to understanding program structure and data flow. Fall back on the same audit methods whenever a decompiler lead dead-ends.

Two habits carry the whole module. First, name the format, the architecture, and the compiler before reading any code, because each fixes what the code can look like and what the mitigations cost. Second, treat every recovered name, type, and prototype as a hypothesis with a confidence, and check the ones the exploit depends on against a run. The reversing deliverable is a written, checkable understanding, not a decompiler window.

\section*{References}
\begin{enumerate}
	\item The Art of Software Security Assessment, McDonald, Schuh \& Dowd (bugs from the auditor's perspective)
	\item Apple's OS X ABI Mach-O File Format Reference (object format ground truth)
	\item https://www.capstone-engine.org (disassembly engine used by many tools)
	\item https://www.jetbrains.com/help/decompiler (and Ghidra: https://ghidra-sre.org)
	\item https://rada.re (radare2), https://rizin.re
	\item The Shellcoder's Handbook, Wiley --- platform specific exploitation basics
	\item Practical Malware Analysis, Sikorski \& Stolz
	\item https://www.fuzzingbook.com (uses reversing to seed fuzzers)
	\item ELF specification (refspecs.linuxfoundation.org), Microsoft PE/COFF specification (docs.microsoft.com/windows/win32/debug/pe-format)
	\item DWARF Debugging Information Format (dwarfstd.org) and the Itanium C++ ABI (itanium-cxx-abi.github.io)
\end{enumerate}

\end{document}
